\documentclass{beamer}

\DeclareMathOperator{\SF}{sf}

\title{Chapter 2 -- Bayes' Rule}
\date{DATA 335 -- Statistical Modeling -- Winter 2026}

\begin{document}

\frame{\titlepage}

\begin{frame}
    \frametitle{Bayes' Rule}
    Let $A$ and $B$ be events. Then:
    \[
    P(B\mid A) = \frac{P(A\mid B)P(B)}{P(A)}
    \]
\end{frame}

\begin{frame}
    \frametitle{The law of total probability}
    Suppose that events $B_1,\ldots,B_n$ partition the sample space $\Omega$:
    \[
    \Omega = \bigcup_{i=1}^n B_i,\qquad B_i\cap B_j=\varnothing\quad(i\neq j).
    \]
    The \emph{law of total probability}:
    \[
    P(A) = \sum_{i=1}^n P(A\mid B_i)P(B_i).
    \]
    
    Substituting for the denominator in Bayes' Rule gives:
    \[
    P(B\mid A) = \frac{P(A\mid B)P(B)}{\displaystyle\sum_{i=1}^n P(A\mid B_i)P(B_i)}.
    \]
\end{frame}

\begin{frame}
    \setlength\parskip{1em}
    \frametitle{Example: Fake news}
    Given:
    \begin{itemize}
        \setlength\itemsep{1em}
        \item 2 news articles in 5 are fake.
        \item 4 in 15 fake news articles have an exclamation point in the title.
        \item 1 in 45 real news articles have an exclamation point in the title.
    \end{itemize}
    What is the probability that an article with an exclamation point
    in the title is fake?
    
\end{frame}


\begin{frame}
    \setlength\parskip{1em}
    Let $A$ be the event \textit{the article contains an exclamation point in its title}.

    Let $B_1$ be the event \textit{the article is fake news}.

    Let $B_2$ be the event \textit{the article is real news}.

    Since an article is either real or fake but not both,
    $$
    B_1\cup B_2 = \Omega\qquad\text{and}\qquad
    B_1\cap B_2 = \varnothing.
    $$
\end{frame}
    
\begin{frame}
    Since $2$ in $5$ news articles are fake,
    $$
    P(B_1) = \frac25\qquad\text{and}\qquad P(B_2) = 1 - \frac25 = \frac35
    $$

    Since $4$ in $15$ (resp., $1$ in $45$) fake (resp., real) articles
    have an exclamation point in the title,
    \[
    P(A\mid B_1) = \frac{4}{15}\qquad\text{and}\qquad
    P(A\mid B_2) = \frac{1}{45}.
    \]
\end{frame}

\begin{frame}
    Apply Bayes' Rule to compute the probability that an
    article with an exclamation point in its title is fake news:
    
    \begin{align*}
    P(B_1\mid A) &= \frac{P(A\mid B_1)P(B_1)}{P(A\mid B_1)P(B_1) + P(A\mid B_2)P(B_2)}\\
    &= \frac{\frac{4}{15}\times \frac{2}{5}}{\frac{4}{15}\times \frac{2}{5} + \frac{1}{45}\times \frac{3}{5}}\\
    &= \frac89.
    \end{align*}

    Thus, $8$ in $9$ articles with exclamation points in their titles are fake news.
\end{frame}


\begin{frame}
    \setlength\parskip{1em}
    \frametitle{Exercise: Vampirism}
    A new test for vampirism has the following characteristics:
    \begin{itemize}
        \setlength\itemsep{1em}
        \item A vampire tests positive for vampirism $95\%$ of the time.
        \item A mortal tests positive for vampirism $1\%$ of the time.
    \end{itemize}
    Given that vampires make up $0.1\%$ of the population, what's the probability that someone who tests positive for vampirism is, in fact, a vampire?
\end{frame}


\begin{frame}    
    \begin{align*}
    &P(\text{vampire}\mid+)\\
    &=\frac{P(+\mid\text{vampire})P(\text{vampire})}
    {P(+\mid\text{vampire})P(\text{vampire}) + P(+\mid\text{mortal})P(\text{mortal})}\\
    &= \frac{0.95\times 0.001}{0.95\times 0.001 + 0.01\times 0.999}\\
    &= 0.087
    \end{align*}

    Thus, an individual that tests positive for vampirism has an $8.7\%$ chance of being a vampire.
    Said differently, $91.3\%$ of positive test results are \emph{false positives}.
\end{frame}

\begin{frame}
    \frametitle{Example: Pop versus Soda}
\end{frame}

\begin{frame}
\begin{itemize}
    \item Men's heights in inches are normally distributed with mean $70$
    and standard deviation $3$.
    \item Women's heights in inches are normally distributed with mean $64.5$
    and standard deviation $2.5$.
\end{itemize}    
Assuming that men and women are equally represented in the population,
what's the probability that a person $>68$ inches tall is a man?
\end{frame}

\begin{frame}
    \setlength\parskip{1em}
    We introduce random variables $X$ for height and $Y$ for sex,
    the latter coded as $0$ for female and $1$ for male.

    Note that the events $Y=0$ and $Y=1$ are mutually exclusive
    and exhaust all possibilities.

    By Bayes' Rule,
    \begin{align*}
    &P(Y=1\mid X>68)\\
    &=\frac{P(X>68\mid Y=1)P(Y=1)}
    {P(X>68\mid Y=0)P(Y=0) + P(X>68\mid Y=1)P(Y=1)}.
    \end{align*}

    Since men and women are equally represented in the population,
    \[
    P(Y=0) = 0.5 = P(Y=1).
    \]
\end{frame}

\begin{frame}
    \setlength\parskip{1em}
    The conditional probabilities $P(X>x\mid Y=y)$ are values of
    \emph{survival functions} of the normal distributions.

    The survival function $\SF(x\mid \mu,\sigma)$ of the
    normal distribution with mean $\mu$ and standard deviation $\sigma$ is:
    \[
    \SF(x\mid \mu, \sigma) = P(X>x),\qquad
    X\sim N(\mu, \sigma)
    \]

    Thus,
    \begin{align*}
    P(X>68\mid Y=0) &= \SF(68\mid 64.5, 2.5),\\
    P(X>68\mid Y=1) &= \SF(68\mid 70, 3).
    \end{align*}

    Compute values of survival functions using \texttt{scipy.stats}.
\end{frame}

\begin{frame}
    \frametitle{Theorem}
    Suppose
    \begin{align*}
    X&\sim N(\mu_0,\sigma_0^2),\\
    Y\mid X=x&\sim N(x, \sigma^2).
    \end{align*}

    Then
    \[
    X\mid Y=y\sim N(M_y, S^2),
    \]
    where
    \[
    M_y = \frac{\dfrac1{\sigma_0^2}\mu_0 + \dfrac1{\sigma^2}y}{\dfrac1{\sigma_0^2} + \dfrac1{\sigma^2}}
    \qquad\text{and}\qquad
    S^2 = \frac{1}{\dfrac1{\sigma_0^2} + \dfrac1{\sigma^2}}.
    \]
\end{frame}

\begin{frame}
    \setlength\parskip{1em}
    Suppose that intrinsic intelligence quotient (IQ) is normally distributed in the
    population with mean $100$ and standard deviation $15$.

    Suppose, further, that IQ test scores for a person with an IQ of $x$
    are normally distributed with mean $x$ and standard deviation $5$.

    What is the probability that a person who scores $120$ on an IQ test has IQ at least 120?
    
\end{frame}

\begin{frame}
    Let $X$ denote intrinsic IQ and let $Y$ denote test score.

    Then
    \begin{align*}
    X&\sim N(100,15^2),\\
    Y\mid X=120&\sim N(120, 5^2).
    \end{align*}

    By the Theorem,
    \[
    X\mid Y=y\sim N(M_y, S^2),
    \]
    where
    \[
    M_y = 118
    \qquad\text{and}\qquad
    S^2 = 22.5 = 4.74^2.
    \]
    It follows that
    \[
    P(X > 120\mid Y=120) = \SF(120\mid 118, 4.74) = 0.34.
    \]
\end{frame}

\end{document}